
\subsubsection{RLHF Evaluation and Benchmarks}

Reinforcement Learning from Human Feedback has become the standard approach for aligning language models with human preferences. Ouyang et al. (2022) introduced InstructGPT using RLHF to improve helpfulness and safety. Evaluation typically relies on aggregate benchmark metrics---TruthfulQA measures factual accuracy \citep{lin2022truthfulqa}, ETHICS assesses moral reasoning \citep{hendrycks2021ethics}, and HHH evaluates helpfulness, harmlessness, and honesty \citep{askell2021general}.

These benchmarks report aggregate accuracy without task-level analysis. The present work clusters tasks by behavioral signals (calibration patterns), enabling mechanism investigation at the task-type level rather than aggregate statistics.

\subsubsection{Model Calibration}

Calibration measures alignment between model confidence and actual correctness. Guo et al. (2017) established modern calibration analysis for neural networks. Prior work treats miscalibration as a general phenomenon. The current study shows calibration inversion clusters systematically (silhouette = 0.6016), indicating specific task types trigger miscalibration rather than uniform overconfidence.

\subsubsection{Bidirectional Alignment Framework}

Shen et al. (2024) proposed a theoretical framework distinguishing AI-to-Human alignment from Human-to-AI alignment based on a survey of 400+ interdisciplinary papers. The present work operationalizes this framework, hypothesizing that tasks requiring bidirectional adaptation correlate with calibration inversion. While keyword-based feature detection proved insufficient, the underlying mechanism chain is verified.

\subsubsection{Positioning}

Prior work either evaluates RLHF outcomes or analyzes mechanisms in isolation. This work bridges these by using behavioral patterns to investigate training mechanisms.
